\underline{Riemann-Stieltjes Integral}: If $f$ is continuous and $g$ is of \underline{finite variation}, then $\int_{a}^{b} f(x) ~ \mathrm{d}g(x)$ is well-defined, and is given by
$$
    \int_{a}^{b} f(x) ~ \mathrm{d}g(x):=\lim_{n\to \infty} \sum_{i=0}^{n-1}f(\xi_{i})\left(g(x_{i+1}) - g(x_{i})\right)
$$
where $a \leq x_{0} < x_{1} < \cdots < x_{n} \leq b$ and $\xi_{i} \in [x_{i}, x_{i+1}]$.

\underline{Main Task}: define the It\^o integral
$$
    (H\cdot X)_{t} = \int_{0}^{t} H_{s} ~\mathrm{d}X_{s}.
$$

The most important case is when $X$ is Brownian motion.

\thmnn{$\displaystyle\sum_{k=1}^{2^{n}}\left(B_{\frac{k}{2^{n}}} - B_{\frac{k-1}{2^{n}}}\right)^{2} \xlongrightarrow{a.s.} 1 .$}

Recall that the finite variation of a function $g:[a,b] \to \mathbb{R}$ is
$$
    V_{g}([a,b]) = \int_{a}^{b} \abs{ \mathrm{d}g(x)} = \sup_{n}\sum_{i=1}^{n}\abs{g(x_{i+1}) - g(x_{i})}
$$

\cornnp{$V_{B}([0,1]) = \infty$ a.s.}{
    \begin{equation*}
        \begin{aligned}
                    \textcolor{red}{(\star)}\sum_{k=1}^{2^{n}}\left( B_{\frac{k}{2^{n}}} - B_{\frac{k-1}{2^{n}}}\right)^{2} &\leq \sup_{1 \leq k \leq 2^{n}} \abs{B_{\frac{k}{2^{n}}} - B_{\frac{k-1}{2^{n}}}}\cdot \sum_{k=1}^{2^{n}}\abs{B_{\frac{k}{2^{n}}} - B_{\frac{k-1}{2^{n}}}}\\
        & \leq V_{B}([0,1]) \cdot \sup_{1 \leq k \leq 2^{n}}\abs{B_{\frac{k}{2^{n}}} - B_{\frac{k-1}{2^{n}}}}
        \end{aligned}
    \end{equation*}
    If $V_{B}([0,1]) < \infty$, then by the continuity of Brownian paths,
    $$
        \sup_{1 \leq k \leq 2^{n}}\abs{B_{\frac{k}{2^{n}}} - B_{\frac{k-1}{2^{n}}}} \to 0,
    $$
    and hence $\textcolor{red}{(\star)} \to 0$.
    
    This contradicts $\textcolor{red}{(\star)} \to 1$. Therefore, with probability one,
    $$
        V_{B}([0,1]) = \infty.
    $$
}

This shows that Brownian motion is too irregular to be treated by the classical Riemann--Stieltjes theory. Therefore, to define
$$
    (H\cdot X)_{t} = \int_{0}^{t} H_{s}~\mathrm{d}X_{s},
$$
we need a new framework, where typically $H$ is previsible (predictable) and $X$ is a semimartingale.

\ex{}{Let $H_{n} \in \mathcal{F}_{n-1}$ and let $(M_{n})$ be a martingale.

Then $(H\cdot M)_{n} = \sum_{k=1}^{n}H_{k}(M_{k} - M_{k-1})$ is a martingale.}

\section{Set-up}

Throughout, let $(\Omega, \mathcal{F}, \mathsf{P})$ be a probability space. A stochastic process $X=(X_{t})_{t \geq 0}$ is a family of random variables taking values in $(\mathbb{R}^{d},\mathcal{B}(\mathbb{R}^{d}))$:
$$
\begin{aligned}
    X\colon [0,\infty) \times \Omega &\to \mathbb{R}^{d}\\
   (t,\omega) &\mapsto X_{t}(\omega).
\end{aligned}
$$

For each fixed $\omega$, the map $t \mapsto X_{t}(\omega)$ is called the sample path of $X$ corresponding to $\omega$.

\defnn{For two stochastic processes $X$ and $Y$, we say that $Y$ is a modification of $X$ if
$$
    \forall~ t \geq 0,\quad \mathsf{P}(X_{t}= Y_{t}) = 1.
$$}

\defnn{We say that $X$ and $Y$ are indistinguishable if
$$
    \mathsf{P}(X_{t} = Y_{t},~ \forall~ t \geq 0) = 1.
$$
In this case, $Y$ is also called a version of $X$.}

\ex{}{Let $Y$ be a modification of $X$, and suppose that both $X$ and $Y$ have a.s. right-continuous sample paths. Then $X$ and $Y$ are indistinguishable.}

\defnn{A function $f$ is right-continuous if
$$
    f(x)= f(x+) = \lim_{t\downarrow x} f(t).
$$}

\pf{Since $\forall~  t \geq 0$, $\mathsf{P}(X_{t} \neq Y_{t}) = 0$, the two processes agree almost surely at every fixed rational time. More precisely,
    $$
        \mathsf{P}\left(\bigcup_{q\in \mathbb{Q}^{+}} X_{q} \neq Y_{q}\right) \leq \sum_{q\in \mathbb{Q}^{+}} \mathsf{P} (X_{q} \neq Y_{q}) = 0.
    $$
    Therefore,
    $$
        \mathsf{P}\left(\bigcap_{q\in \mathbb{Q}^{+}} X_{q} = Y_{q}\right) = 1.
    $$
    On the other hand, if two right-continuous functions $f$ and $g$ agree on $\mathbb{Q}^{+}$, then they must agree on all of $[0,\infty)$, since the rationals are dense.

    Applying this observation pathwise to the sample paths of $X$ and $Y$, we conclude that
    $$
        \mathsf{P}\left(X_{t} = Y_{t},~ \forall~ t \geq 0\right) = 1.
    $$
}

\ex{}{Let $T \sim \mathrm{Exp}(1)$. Define 
    $$
    \begin{aligned}
        &X_{t} \equiv 0,\quad\forall~ t \geq 0 \\
        &Y_{t} = \begin{cases}
            0,& t \neq T\\
            1,&t = T.
        \end{cases}
    \end{aligned}
    $$
    Then for every fixed $t \geq 0$,
    $$
        \mathsf{P}(X_{t} \neq Y_{t}) = \mathsf{P}(Y_{t} = 1) = \mathsf{P}(t = T) = 0.
    $$
    Hence $Y$ is a modification of $X$. However, $X$ and $Y$ are not indistinguishable, because on the event $\{T<\infty\}$ they differ at time $T$.
}
